% Bot ladder -- generated tables/figures under results/elo/ and results/case_studies/

\subsection{Protocol}
\label{sec:elo-protocol}

Full-game strength is the ground truth that every other section proxies. All
registered models compete except \code{v2\_knapsack} and \code{v3}, which are
code-identical to \code{v4}. Each game is a three-player \emph{free-for-all}
on a fresh Poisson map: one seat runs model A (labeled ``solo'') and the other
two seats run \emph{independent} instances of model B --- not a team; the
copies share no state, do not coordinate, and fight each other as readily as
they fight A. Model A scores 1 if it holds strictly more nodes at game end
(elimination or 5000 ticks) than the \emph{better-performing} of the two B
copies, $\tfrac12$ on a tie. Duplicating the opponent gives every game an
internal baseline (copy vs.\ copy on the same map) and makes the environment
majority-hostile to the rated model, which separates models faster than a
symmetric duel at these game counts. Each unordered pair plays four games,
half with each model in the solo seat, with deterministic seeds; the full log
is \code{results/elo/raw/games.csv}.

Ratings are fit by maximum likelihood under a Bradley--Terry model
\citep{bradleyterry1952} via minorization--maximization \citep{hunter2004} on
the aggregate score matrix, then mapped onto the Elo scale
($1500 + 400\log_{10}p$) \citep{elo1978}. Unlike sequential K-factor updates,
the fit is independent of game order and exactly reproducible from the log.

\subsection{The ladder}
\label{sec:elo-ladder}

\begin{table}[ht]
\centering
\resultstable{elo/tables/elo_table}
\caption{Bradley--Terry ratings on the Elo scale. ``Score rate'' is mean score
per game across all opponents. Read in tiers: with four games per pair, gaps
under $\sim$100 Elo are noise, and the extreme negative values at the bottom
are scale artifacts of near-zero score rates.}
\label{tab:elo}
\end{table}

Four headlines:

\begin{enumerate}
  \item \textbf{\code{v8} tops the ladder} (score rate 0.90) --- and \code{v8}
    is the \emph{only} model using the distance-decayed local softmax, a
    branch abandoned when v9 pivoted to OT transport. The most valuable single
    tweak tested is currently unused by every later model.
  \item \textbf{OT transport is the strongest lineage trunk}: \code{v9},
    \code{v12}, \code{v10} fill ranks 2--5 and clear every pre-OT model except
    \code{v8}.
  \item \textbf{\code{v11} is a regression, not noise} (0.17, adrift from
    parent \code{v10} at 0.65 and child \code{v12} at 0.75) --- dissected in
    \S\ref{sec:elo-v11}.
  \item \textbf{Pacifism is underrated at these game counts}:
    \code{v0\_expansion}, which has no war agent at all, sits mid-table above
    \code{v4}--\code{v6}; the knapsack-era aggressors are the tournament's
    losers.
\end{enumerate}

\subsection{How transitive is the game? (HodgeRank)}
\label{sec:elo-hodge}

A scalar rating assumes the tournament is a gradient: everyone's strength
lies on one axis. The pairwise data can measure that assumption. Take the
skew-symmetric edge flow $Y_{ij} = (\text{score rate of } i \text{ over }
j) - \tfrac12$ on the complete comparison graph and Hodge-decompose it
\citep{jiang2011}: a gradient part (explained by a potential $s_i$, i.e.\ by
\emph{some} ranking) plus a residual, which on a complete graph is pure curl
--- local rock-paper-scissors cycles that no ranking can represent.

The result: the gradient explains \textbf{\hodgetransitivity} of the flow's
energy; \textbf{\hodgecurl{} of the tournament is cyclic}. The Hodge potential
and the Bradley--Terry fit agree exactly on order (Spearman
$\rho = \hodgespearman$; Figure~\ref{fig:hodge}), so Elo is not mis-ranking
anyone --- the \hodgecurl{} is simply invisible to \emph{any} scalar rating,
a property of the game, not of the estimator.

\begin{figure}[ht]
\centering
\includegraphics[width=0.62\textwidth]{\resultsfig{elo/figures/hodge_vs_elo.pdf}}
\caption{Hodge potential vs.\ Bradley--Terry rating. The monotone relationship
(Spearman $\rho = \hodgespearman$) says the two scalar summaries agree; the
interesting information is what neither can hold --- the cyclic
\hodgecurl{} quantified by the decomposition.}
\label{fig:hodge}
\end{figure}

\begin{table}[ht]
\centering
\resultstable{elo/tables/curl_triangles}
\caption{Largest-curl triangles (residual flow around the 3-cycle, in
score-rate units; $\pm\tfrac32$ is the maximum possible). The podium itself
cycles: \code{v12} beats \code{v8} beats \code{v9} beats \code{v12}.}
\label{tab:curl}
\end{table}

\subsection{Case study: the podium cycle}
\label{sec:elo-podium}

The top three form one of the largest cycles in the tournament
(Table~\ref{tab:curl}): head-to-head, \code{v9} took its pair against
\code{v8} 0.75, despite \code{v8} out-rating it against the field.
Figure~\ref{fig:v9v8} replays tournament game 1272 (\code{v8} solo vs.\ two
\code{v9}). The game is decided almost immediately: \code{v8} expands fastest
off the start --- peaking near 75 nodes by tick 250 --- but is torn apart in
the first contact wars and eliminated by tick $\sim$600. Everything after is
the two \code{v9} copies partitioning the map (the solid copy plateaus around
240 nodes and wins the node count at timeout). The qualitative lesson: losses
in this protocol are not slow grinds but early-war collapses, so whatever
edge \code{v8}'s local softmax buys against the field, it does not survive
simultaneous pressure from two OT-supplied fronts on this map.

\begin{figure}[ht]
\centering
\includegraphics[width=\textwidth]{\resultsfig{case_studies/figures/v9_beats_v8.pdf}}
\caption{Tournament game 1272 replayed. All three curves are the three
players of this \emph{single} game: \code{v8} in the solo seat (blue) and the
two independent \code{v9} instances (green, solid/dashed). Left: total
troops; right: nodes held. After \code{v8}'s elimination the remaining
dynamics are the two \code{v9} seats fighting each other.}
\label{fig:v9v8}
\end{figure}

\subsection{Case study: the \code{v11} regression, with counterfactual}
\label{sec:elo-v11}

\code{v11} differs from \code{v10} by exactly one feature: OT demands are
netted against \emph{effective} (in-transit-inclusive) troop counts.
Figure~\ref{fig:v11} replays tournament game 1301, which \code{v11} lost, and
then replays the \emph{same seed, same map, same slot, same opponents} with
\code{v12} substituted --- a single-variable comparison, since \code{v12} is
\code{v11} plus a standing frontline-garrison demand.

The contrast is stark. In the tournament game, \code{v11}'s trajectory is
nearly indistinguishable from the \emph{losing} \code{v10} copy --- both
expand to $\sim$100 nodes by tick 300 and are wiped out by tick 1000 while
the surviving \code{v10} sweeps the map. \code{v11} does not exhibit a
visibly pathological curve; it simply fights like a mid-tier bot and loses.
But substitute \code{v12} into the identical position and the same game is a
rout \emph{in the other direction}: all 347 nodes by tick 1000, both
\code{v10} copies dead by tick 900. One added demand term converts this map
from a loss into total domination, and in aggregate from a 0.17 score rate to
0.75. The mechanism this implicates --- with netting alone, arriving
reinforcements cancel the very demands that summoned them, so contested nodes
stop pulling troops exactly when they need them, while \code{v12}'s standing
garrison demand keeps the pull alive --- fits both the aggregate and the
counterfactual, but remains a hypothesis until the OT demands are instrumented
directly in a replayed game (\S\ref{sec:open}, item 2).

\begin{figure}[ht]
\centering
\includegraphics[width=\textwidth]{\resultsfig{case_studies/figures/v11_stall.pdf}}\\[2pt]
\includegraphics[width=\textwidth]{\resultsfig{case_studies/figures/v12_counterfactual.pdf}}
\caption{Controlled counterfactual on tournament seed 1301: two separate
games, one per row, each row showing all three players of that game. Top:
the tournament game --- \code{v11} in the solo seat (blue) vs.\ two
independent \code{v10} instances (green); \code{v11} and one \code{v10} are
both destroyed by the other \code{v10}. Bottom: a replay with \code{v12}
substituted into the identical seat --- a rout in the other direction. The
only delta between the two solo models is the frontline-garrison demand.}
\label{fig:v11}
\end{figure}
